\chapter{Uncertainty Quantification}



This chapter will investigate the fundamentals of uncertainty quantification. Techniques and methods relevant to the
prediction of DP performance will be identified and documented.

\section{What is Uncertainty Quantification?}
A typical uncertainty quantification problem involves a mathematical model, subject to
uncertainty about the correct parameter values \parencite{sullivanIntroductionUncertaintyQuantification2015}.
\textcite{higdonUncertaintyQuantificationError2009} define uncertainty quantification as follows:

\begin{quotation}
    "Uncertainty quantification (UQ) studies all sources of error and uncertainty, including the following: systematic
    and stochastic measurement error; ignorance; limitations of theoretical models; limitations of numerical
    representations of those models; limitations of the accuracy and reliability of computations, approximations, and
    algorithms; and human error. A more precise definition is UQ is the end-to-end study of the reliability of
    scientific inferences"
    \parencite{higdonUncertaintyQuantificationError2009}.
\end{quotation}

In general in the study of UQ, one is interested in the relationship between the input and output. Two general examples
are the forward and the inverse problem \parencite{sullivanIntroductionUncertaintyQuantification2015}. In both examples
we will consider a general model of a physical system:
\[
z = f(x, y)
\]

In the forward problem, $f, \, x, \, y$ are known. $z$ can be calculated. However, if $f$ is an imperfect model or $x,
\, y$ are subject to some uncertainty, then this uncertainty will in some way affect $z$. This is known as the 
\textit{forward propegation of uncertainty}.

In the inverse problem $z, \, x,\, y$ are measured. We are interested in determining $f$. These measurements are subject to noise and as $f$ is a model of a
real system and thus has infinitely many degrees of freedom, the problem is undetermined. 


\section{Sensitivity analysis}
In this section sensitivity analysis will be discussed. Sensitivity analysis is the understanding of how a function
$f(x_1, \dots, x_n)$ depends on variations in $x_i$, both individually and in combinanation with other $x_i$
\parencite{sullivanIntroductionUncertaintyQuantification2015}. There are two approaches: local sensitivity analysis and
global sensitivity analysis. Local analyses study how the system responds to input variations around a point
$\overbar{x}$. Global sensitivity analyes study the average sensitivity of $f$ to variations across the entire domain. A
practical application is that understanding the sensitivity of a system can allow an engineer to decide which variables
to include uncertainty of in their uncertainty modelling. 

\section{Intrusive Methods}
Intrusive methods for uncertainty quantification are methods that require alteration of the governing equations ot
incorporate uncertainty directly.

\section{Polynomical Chaos Expansions}
Polynomial chaos expansions (PCE) are a spectral method for uncertainty quantification. A spectral mathod is a method where a
stochastic process is represented as a series expansion using an orthogonal basis of functions. Spectral methods
decompose a random variable into components, similar to a fourier analysis \parencite{heuvelineHybridGeneralizedPolynomial2014}.

PCE uses polynomials as the orthogonal basis

\section{Non-intrusive methods}
Non-intrusive methods for uncertainty quantification are methods that allow the use of existing models as a black box
without adjusting their internal structure for UQ. 

\textcite{sullivanIntroductionUncertaintyQuantification2015} discusses several of these methods. The guiding principle
of these is that the function is evaluated at a number of points. A surrogate model is then made based on these samples.
The first method discussed is a non-intrusive spectral projection. With this method the aim is to select the best
polynomial chaos coefficients so as to model the function best using a set of orthogonal polynomials.
The second method, stochastic collocation, creates an interpolant between chosen nodes. The third method is gaussian
process regression. This method, in contrast to the two previous methods, does not assume perfectly accurate
observations. The function is assumed to be a random gaussian field and any observations made reduce the uncertainty at
those points.

One important aspect to note is that these methods all assume the function does not change in time. This makes these
methods as is unsuitable for dynamic systems. 

